\documentclass[11pt]{ctexart}
\usepackage{amsmath,amssymb,amsthm}
\usepackage[margin=2.5cm]{geometry}
\usepackage{booktabs}

\theoremstyle{plain}
\newtheorem{theorem}{Theorem}[section]
\newtheorem{lemma}[theorem]{Lemma}
\newtheorem{proposition}[theorem]{Proposition}
\newtheorem{corollary}[theorem]{Corollary}
\newtheorem{conjecture}[theorem]{Conjecture}
\theoremstyle{definition}
\newtheorem{definition}[theorem]{Definition}
\newtheorem{example}[theorem]{Example}
\newtheorem{openproblem}[theorem]{Open problem}
\theoremstyle{remark}
\newtheorem{remark}[theorem]{Remark}

\newcommand{\kk}{\mathsf{k}}
\newcommand{\gl}{\mathfrak{g}}
\newcommand{\hl}{\mathfrak{h}}
\newcommand{\LL}{\mathfrak{L}}
\newcommand{\UU}{\mathcal{U}}
\newcommand{\Sym}{\operatorname{Sym}}
\newcommand{\Der}{\operatorname{Der}}
\newcommand{\GKdim}{\operatorname{GKdim}}

\title{Noetherianity of universal enveloping algebras:\\[2pt]
\large finite dimension, finite codimension, and graded obstructions}
\author{G}
\date{}

\begin{document}
\maketitle

\begin{abstract}
Let $\kk$ be a field of characteristic zero and $\gl$ a Lie algebra over $\kk$ with
universal enveloping algebra $\UU(\gl)$. This note is about the following question:
does $\UU(\gl)$ noetherian force $\dim_\kk\gl<\infty$? We give self-contained,
line-by-line proofs of four things: the classical positive direction (finite
dimension implies noetherian, through the Poincar\'e--Birkhoff--Witt filtration,
the Hilbert basis theorem and a lifting lemma for filtered rings); the abelian
counterexample (infinite dimension implies not noetherian, with an explicit
strictly ascending chain); a descent lemma for ``large'' ring extensions; and the
module-theoretic facts around finite-codimension subalgebras, together with an
example showing why the naive module-finiteness route fails. We then state, under
the hypotheses under which they are published and with an explicit note that their
proofs are not reconstructed here, the graded obstruction for simple graded Lie
algebras with finite-dimensional homogeneous components, and its special cases
(Krichever--Novikov algebras, graded algebras of polynomial growth, finitely
generated infinite-step solvable algebras). The general converse,
``$\UU(\gl)$ noetherian implies $\dim_\kk\gl<\infty$'', is stated as a conjecture
and is \emph{not} proved or disproved here; we say exactly which step is missing.
All proofs in this note are natural-language arguments written and reviewed line by
line by hand.
\end{abstract}

\section{Introduction}

Let $\kk$ be a field of characteristic $0$, let $\gl$ be a Lie algebra over $\kk$ and
let $\UU(\gl)$ be its universal enveloping algebra. Recall that a ring is
\emph{noetherian} if its left ideals satisfy the ascending chain condition; for
$\UU(\gl)$ the left and the right condition are equivalent, because the principal
anti-automorphism $x\mapsto -x$ ($x\in\gl$) exchanges left and right ideals.

The question motivating this note is whether noetherianity of $\UU(\gl)$ detects
finite dimension:
\begin{equation}\label{eq:Q}
\UU(\gl)\ \text{noetherian}\ \Longrightarrow\ \dim_\kk\gl<\infty\ ?
\end{equation}
Two extreme cases have been known for a long time. If $\dim_\kk\gl=n<\infty$, then
$\UU(\gl)$ is noetherian; this is the standard route through the
Poincar\'e--Birkhoff--Witt filtration and the Hilbert basis theorem, and it is
recorded for instance in \cite[Chapter 1]{MR01}. If $\gl$ is abelian, then
$\UU(\gl)\cong\Sym(\gl)$ is a polynomial algebra in $\dim_\kk\gl$ variables, which is
noetherian exactly when the number of variables is finite. The interesting case is
therefore: can an infinite-dimensional \emph{non-abelian} $\gl$ have a noetherian
enveloping algebra?

The known negative results are obstruction results of a graded nature. In
\cite{SW14} it is shown that the universal enveloping algebra of the Witt algebra
$\Der\kk[t,t^{-1}]$ is not noetherian, by exhibiting a strictly ascending sequence
of two-sided ideals; the authors also formulate the conjecture that no
infinite-dimensional Lie algebra over a field of characteristic zero has a
noetherian enveloping algebra. Generalizations cover simple Lie algebras graded by
$\mathbb{Z}^n$ with finite-dimensional homogeneous components, the
Krichever--Novikov algebras of \cite{KN87}, graded algebras of polynomial growth,
and finitely generated infinite-step solvable algebras. The common feature is
that noetherianity is refuted by an explicit chain that is visible in the grading.

What this note does and does not do. It gives complete proofs of the positive
direction (Section 3, Conjectures~\ref{thm:findim} and~\ref{thm:abelian}), of two
general ring-theoretic lemmas used everywhere (Lemmas~\ref{lem:hilbert}
and~\ref{lem:lift}), and of the module-theoretic facts surrounding
finite-codimension subalgebras (Proposition~\ref{prop:codim}), including an example
(Example~\ref{ex:nomodule}) showing that a finite-codimension subalgebra does
\emph{not} in general make the bigger enveloping algebra a finitely generated module
over the smaller one. It then reports the graded obstruction
(Conjecture~\ref{thm:gradedsimple}) under its published hypotheses with an explicit
statement that the detailed argument is not reconstructed here, and derives from it
the Witt case (Corollary~\ref{cor:witt}) and lists the three classes mentioned above
(Conjectures~\ref{thm:kn}, \ref{thm:poly}, \ref{thm:solvable}). The general question
\eqref{eq:Q} is left open and is stated as Conjecture~\ref{conj:SW} together with an
explicit list of what is missing (Section~5).

Numerical evidence for the two finite checks referred to below is collected in
Appendix~A; any conclusion resting on it depends on that evidence and on the stated
bounds, and nothing in Appendix~A is used in the proofs of Section~3.

\section{Definitions and notation}\label{sec:defs}

\begin{definition}[Filtered and graded rings]\label{def:filtered}
A \emph{filtered $\kk$-algebra} is a $\kk$-algebra $A$ together with $\kk$-subspaces
$A_0\subseteq A_1\subseteq A_2\subseteq\cdots$ such that $1\in A_0$,
$A_iA_j\subseteq A_{i+j}$ and $\bigcup_i A_i=A$. Its \emph{associated graded algebra}
is $\operatorname{gr}A=\bigoplus_{n\ge0}A_n/A_{n-1}$ (with $A_{-1}=0$), with the
product induced by $A_i\times A_j\to A_{i+j}$. A \emph{graded $\kk$-algebra} is a
$\kk$-algebra $G=\bigoplus_{n\ge0}G_n$ with $G_iG_j\subseteq G_{i+j}$.
A left ideal $I$ of a filtered algebra $A$ carries the induced filtration
$I_n=I\cap A_n$ and one writes $\operatorname{gr}I=\bigoplus_n (I_n+A_{n-1})/A_{n-1}$,
a graded subspace of $\operatorname{gr}A$.
\end{definition}

\begin{definition}[Noetherian]\label{def:noetherian}
A ring $R$ is \emph{left noetherian} if every ascending chain
$I_1\subseteq I_2\subseteq\cdots$ of left ideals of $R$ is eventually constant;
equivalently, if every left ideal of $R$ is finitely generated. Right noetherian is
defined symmetrically. For a commutative ring the two notions coincide and one
speaks simply of a noetherian ring.
\end{definition}

\begin{definition}[Enveloping algebra]\label{def:U}
For a Lie algebra $\gl$ over $\kk$, $\UU(\gl)$ is the quotient of the tensor algebra
$T(\gl)$ by the two-sided ideal generated by $xy-yx-[x,y]$, $x,y\in\gl$. We use the
\emph{Poincar\'e--Birkhoff--Witt (PBW) theorem} in the following form: if
$(x_i)_{i\in I}$ is a totally ordered $\kk$-basis of $\gl$, the ordered monomials
$x_{i_1}^{a_1}\cdots x_{i_r}^{a_r}$ with $i_1<\cdots<i_r$ and $a_j\ge1$, together
with $1$, form a $\kk$-basis of $\UU(\gl)$. The \emph{standard filtration} is
$\UU_n(\gl)=\kk$-span of all products of at most $n$ elements of $\gl$.
\end{definition}

\begin{definition}[Witt algebras, graded Lie algebras]\label{def:witt}
Put $W=\Der\kk[t,t^{-1}]$, the Lie algebra of derivations of the Laurent polynomial
ring, with basis $d_m=t^{m+1}\partial_t$, $m\in\mathbb{Z}$, and bracket
$[d_m,d_n]=(n-m)d_{m+n}$; write $W_{\ge-1}=\Der\kk[t]=\bigoplus_{m\ge-1}\kk d_m$, the
subalgebra of polynomial vector fields. Note that $W_{\ge-1}$ has codimension $1$ in
$W$ inside the span of negative modes, and that both are infinite dimensional.
A Lie algebra $\LL$ is \emph{$\mathbb{Z}^n$-graded} if
$\LL=\bigoplus_{m\in\mathbb{Z}^n}\LL_m$ as vector spaces with
$[\LL_m,\LL_{m'}]\subseteq\LL_{m+m'}$; it is \emph{graded simple} if there is no
two-sided ideal $0\neq J\neq\LL$ with $J=\bigoplus_m (J\cap\LL_m)$, and
$\dim_\kk\LL_m<\infty$ for all $m$ is the ``finite-dimensional homogeneous
components'' hypothesis used below. The Witt algebra is $\mathbb{Z}$-graded with
$\dim_\kk W_m=1$ and is graded simple.
\end{definition}

\begin{remark}[What ``proof'' means in this note]\label{rem:standard}
Every result stated below with a proof has a written, line-by-line argument: each
step is either a computation, a citation of one of the lemmas proved here, or an
explicitly named standard input (PBW, the Hilbert basis theorem, faithful flatness).
Results that are reported from the literature rather than proved here are marked in
their statement by ``under its published hypotheses'' and carry a footnote saying so.
Every proof written out here is a hand-written, line-by-line argument.
\end{remark}

\section{Main results}\label{sec:main}

\subsection{The positive direction}

\begin{lemma}[PBW and the associated graded]\label{lem:pbw}
Let $\gl$ be any Lie algebra over $\kk$ and let $\UU_n=\UU_n(\gl)$ be the standard
filtration of Definition~\ref{def:U}. Then $\bigcup_n\UU_n=\UU(\gl)$,
$\UU_i\UU_j\subseteq\UU_{i+j}$, and $\operatorname{gr}\UU(\gl)\cong\Sym(\gl)$ as
graded $\kk$-algebras. In particular, if $\dim_\kk\gl=n<\infty$ then
$\operatorname{gr}\UU(\gl)\cong\kk[x_1,\dots,x_n]$.
\end{lemma}

\begin{proof}
\emph{(1) Exhaustiveness.} By PBW every element of $\UU(\gl)$ is a finite sum of
ordered monomials of some total degree $d$, and an ordered monomial of total degree
$d$ is a product of $d$ elements of $\gl$, hence lies in $\UU_d$; so
$\bigcup_n\UU_n=\UU(\gl)$. The inclusion $\UU_i\UU_j\subseteq\UU_{i+j}$ holds because
a product of $i$ and $j$ factors is a product of $i+j$ factors.

\emph{(2) The symbol of a monomial depends only on its symmetric product.} For
$x,y\in\gl$ we have $xy-yx=[x,y]\in\gl=\UU_1$, so $xy$ and $yx$ have the same image
in $\UU_2/\UU_1$. Consequently, for any permutation $\sigma$ of $d$ factors,
$\prod x_k$ and $\prod x_{\sigma(k)}$ differ by an element of $\UU_{d-1}$: moving one
factor past another changes the product by a bracket, which lies in $\gl$ and hence
lowers the total degree by at least one. Hence the class of
$x_1\cdots x_d$ in $\UU_d/\UU_{d-1}$ depends only on the symmetric tensor
$x_1\cdots x_d\in\Sym^d(\gl)$, and the resulting map
$\Sym^d(\gl)\to\UU_d/\UU_{d-1}$ is well defined and $\kk$-linear.

\emph{(3) It is an isomorphism.} Surjectivity: $\UU_d$ is spanned by products of $d$
elements of $\gl$, so the images of the monomials span $\UU_d/\UU_{d-1}$.
Injectivity: by PBW the ordered monomials of degree exactly $d$ are linearly
independent in $\UU_d$, and they are not in $\UU_{d-1}$ (their PBW expression has no
terms of smaller total degree), so their classes in $\UU_d/\UU_{d-1}$ are linearly
independent; there are as many of them as there are degree-$d$ monomials in a basis
of $\gl$, i.e.\ $\dim\Sym^d(\gl)$ of them.

\emph{(4) Multiplicativity.} If $a=\prod_{i\le p}x_i$ and $b=\prod_{j\le q}y_j$, then
$ab$ is a product of $p+q$ elements of $\gl$, so by (2) its class equals
$\operatorname{sym}(a)\operatorname{sym}(b)$. Hence $\operatorname{gr}\UU(\gl)$ and
$\Sym(\gl)$ are isomorphic as graded algebras. The last assertion is
$\Sym(\gl)\cong\kk[x_1,\dots,x_n]$ for a basis $x_1,\dots,x_n$ of $\gl$.
\end{proof}

\begin{lemma}[Hilbert basis theorem]\label{lem:hilbert}
If $R$ is a commutative noetherian ring then $R[x]$ is noetherian. In particular,
for a field $\kk$ the polynomial ring $\kk[x_1,\dots,x_n]$ is noetherian for every
$n\ge0$.
\end{lemma}

\begin{proof}
\emph{(0) A field is noetherian.} The only ideals of a field are $0$ and $\kk$, so
every ascending chain stabilises immediately.

\emph{(1) Reduction to a generating set.} Let $I\subseteq R[x]$ be an ideal. Let $L$
be the set of leading coefficients of elements of $I$, together with $0$; $L$ is an
ideal of $R$ (if $a,b\in L$ are leading coefficients of $f,g$, then $a+b$ is the
leading coefficient of $x^{\deg g}f+x^{\deg f}g$, and $ra$ of $rf$). Since $R$ is
noetherian, $L=(a_1,\dots,a_s)$; choose $f_i\in I$ with leading coefficient $a_i$
and degree $d_i$, and put $d=\max_i d_i$. Let
$M=I\cap\{h\in R[x]:\deg h<d\}$; as an $R$-module $M$ is a submodule of the finitely
generated free module $\bigoplus_{k<d}Rx^k$, hence finitely generated (a noetherian
ring is noetherian as a module over itself, finite direct sums of noetherian modules
are noetherian, and submodules of noetherian modules are noetherian). Pick
generators $g_1,\dots,g_t$ of $M$.

\emph{(2) The claim.} We show $I=(f_1,\dots,f_s,g_1,\dots,g_t)$. Let $f\in I$; we
prove $f\in J:=(f_1,\dots,f_s,g_1,\dots,g_t)$ by induction on $m=\deg f$. If
$m<d$ then $f\in M\subseteq J$. If $m\ge d$, let $a$ be the leading coefficient of
$f$. Then $a\in L$, so $a=\sum_i c_i a_i$ with $c_i\in R$, and
$f-\sum_i c_i x^{m-d_i}f_i\in I$ has degree $<m$; by induction it lies in $J$, hence
$f\in J$.

\emph{(3) Conclusion.} Every ideal of $R[x]$ is finitely generated, which for a
commutative ring is equivalent to the ascending chain condition. The statement for
$\kk[x_1,\dots,x_n]$ follows by iterating (0) and (1).
\end{proof}

\begin{lemma}[Lifting noetherianity along a filtration]\label{lem:lift}
Let $A$ be a filtered $\kk$-algebra with exhaustive filtration as in
Definition~\ref{def:filtered}. If $\operatorname{gr}A$ is left noetherian, then $A$
is left noetherian.
\end{lemma}

\begin{proof}
\emph{(1) $\operatorname{gr}I$ is a left ideal of $\operatorname{gr}A$.} Let
$I\subseteq A$ be a left ideal, $I_n=I\cap A_n$. For homogeneous $a\in A_d$ and
$x\in I_n$ we have $ax\in I$ and $ax\in A_{n+d}$, so $ax\in I_{n+d}$; moreover the
class of $ax$ modulo $A_{n+d-1}$ is the product of the classes of $a$ and of $x$
(the products $A_d\times A_n\to A_{d+n}$ are bilinear and compatible with the
quotients). Hence $\operatorname{gr}A\cdot\operatorname{gr}I\subseteq
\operatorname{gr}I$.

\emph{(2) Choosing lifts.} Since $\operatorname{gr}A$ is left noetherian, the graded
left ideal $\operatorname{gr}I$ is generated by finitely many homogeneous elements
$\bar y_1,\dots,\bar y_s$ of degrees $d_1,\dots,d_s$; a generating set of a graded
module may be taken homogeneous because each homogeneous component of
$\operatorname{gr}I$ is spanned by homogeneous pieces of a generating set. Choose
lifts $y_i\in I_{d_i}$.

\emph{(3) The induction.} We prove $I_n=\sum_i A_{n-d_i}y_i$ for all $n\ge0$ (with
$A_m=0$ for $m<0$) by induction on $n$. Let $x\in I_n$. Its class $\bar x\in
(\operatorname{gr}I)_n$ is a sum $\sum_i\bar a_i\bar y_i$ with $\bar a_i$ homogeneous
of degree $n-d_i$ (no generator of degree $>n$ can occur, since all terms have degree
$n$). Choose $a_i\in A_{n-d_i}$ lifting $\bar a_i$. Then
$x'=\sum_i a_iy_i\in I_n$ and $\bar x=\overline{x'}$, i.e.\ $x-x'\in I_{n-1}$. By
induction $x'\in\sum_i A_{n-d_i}y_i$ and $x-x'\in\sum_i A_{n-1-d_i}y_i\subseteq
\sum_i A_{n-d_i}y_i$. Hence $x\in\sum_i A_{n-d_i}y_i$.

\emph{(4) Conclusion.} With $n$ large enough that $x\in A_n$, step (3) gives
$I=\sum_i Ay_i$: $I$ is finitely generated. Since this holds for every left ideal,
$A$ is left noetherian.
\end{proof}

\begin{theorem}[Finite dimension implies noetherian]\label{thm:findim}
Let $\gl$ be a Lie algebra over a field $\kk$ with $\dim_\kk\gl=n<\infty$. Then
$\UU(\gl)$ is left and right noetherian.
\end{theorem}

\begin{proof}
\emph{(1)} By Lemma~\ref{lem:pbw}, $\operatorname{gr}\UU(\gl)\cong\kk[x_1,\dots,x_n]$
with the standard filtration.

\emph{(2)} By Lemma~\ref{lem:hilbert} the ring $\kk[x_1,\dots,x_n]$ is noetherian, so
$\operatorname{gr}\UU(\gl)$ is left noetherian.

\emph{(3)} By Lemma~\ref{lem:lift} applied to $A=\UU(\gl)$, the algebra $\UU(\gl)$ is
left noetherian.

\emph{(4)} The principal anti-automorphism $x\mapsto-x$ of $\gl$ extends to an
anti-automorphism of $\UU(\gl)$ and exchanges left and right ideals; hence
$\UU(\gl)$ is also right noetherian.
\end{proof}

\subsection{The abelian counterexample}

\begin{lemma}[Abelian Lie algebras]\label{lem:abelian}
If $\gl$ is abelian then $\UU(\gl)\cong\Sym(\gl)$ as $\kk$-algebras.
\end{lemma}

\begin{proof}
By PBW the ordered monomials in a basis of $\gl$ form a basis of $\UU(\gl)$. If
$[x,y]=0$ for all $x,y\in\gl$, then two ordered monomials multiply to an ordered
monomial (the factors commute, so reordering changes nothing and no bracket terms
appear), hence the $\kk$-linear map $\Sym(\gl)\to\UU(\gl)$ sending
$x_{i_1}\cdots x_{i_d}$ to the product is a well-defined algebra homomorphism. It is
surjective because the images of all monomials span $\UU(\gl)$ (PBW), and injective
because it sends a basis of $\Sym(\gl)$ to a basis of $\UU(\gl)$.
\end{proof}

\begin{theorem}[Infinite dimension implies not noetherian, abelian case]\label{thm:abelian}
Let $\gl$ be an infinite-dimensional abelian Lie algebra over $\kk$. Then
$\UU(\gl)$ is neither left nor right noetherian.
\end{theorem}

\begin{proof}
\emph{(1)} By Lemma~\ref{lem:abelian} we may identify $\UU(\gl)$ with the polynomial
algebra $\kk[x_i:i\in I]$ in the indeterminates of a basis of $\gl$, where $I$ is
infinite.

\emph{(2)} Since $I$ is infinite, it contains a countably infinite subset; relabel
indeterminates so that $x_1,x_2,x_3,\dots$ are distinct elements of $I$. Put
$I_n=(x_1,\dots,x_n)$, the ideal generated by the first $n$ indeterminates. Then
$I_1\subseteq I_2\subseteq I_3\subseteq\cdots$ is an ascending chain of ideals.

\emph{(3) The chain is strictly ascending.} Fix $n\ge1$. Because $I$ is a basis of
$\gl$, there is a $\kk$-algebra homomorphism $\varphi$ from
$\kk[x_i:i\in I]$ to itself with $\varphi(x_i)=0$ for $i\in\{1,\dots,n\}$ and
$\varphi(x_i)=x_i$ otherwise (a homomorphism out of a polynomial ring is determined
by arbitrary images). Then $\varphi(I_n)=0$ while $\varphi(x_{n+1})=x_{n+1}\neq0$;
hence $x_{n+1}\notin I_n$, i.e.\ $I_n\subsetneq I_{n+1}$.

\emph{(4)} An infinite strictly ascending chain of ideals refutes the ascending chain
condition on both sides (ideals of a commutative ring are left and right ideals).
\end{proof}

\subsection{Finite-codimension subalgebras}

\begin{lemma}[Module finiteness transfers noetherianity]\label{lem:module}
Let $B\subseteq A$ be rings and suppose $A$ is finitely generated as a left
$B$-module. If $B$ is left noetherian, then $A$ is left noetherian.
\end{lemma}

\begin{proof}
If $A=B a_1+\cdots+B a_s$ then $A$ is a quotient of the free left $B$-module $B^s$,
which is noetherian as a finite direct sum of copies of the noetherian module $B$;
quotients of noetherian modules are noetherian (submodules of a quotient correspond
to submodules of the module). A left ideal of $A$ is in particular a left
$B$-submodule of $A$, so any ascending chain of left ideals of $A$ is an ascending
chain of $B$-submodules of the noetherian module $A$ and stabilises.
\end{proof}

\begin{theorem}[Descent along free extensions]\label{thm:descent}
Let $B\subseteq A$ be rings such that $A$ is free as a left $B$-module (equivalently
free as a right $B$-module, in the case at hand). Then $A$ is faithfully flat as a
$B$-module, and for every two-sided ideal $\mathfrak{a}$ of $B$ one has
$A\mathfrak{a}\cap B=\mathfrak{a}$. Consequently, if $A$ is noetherian (i.e.\ has the
ascending chain condition on two-sided ideals) then so is $B$.
\end{theorem}

\begin{proof}
\emph{(1) Faithful flatness.} A free module is flat; it is faithfully flat because
$N\otimes_B A=0$ for a free module of positive rank forces $N=0$ (a free module of
rank $\ge1$ has a direct summand isomorphic to $B$, so $N$ embeds in $N\otimes_B A$).

\emph{(2) The contraction formula.} Let $\mathfrak{a}\subseteq B$ be a two-sided
ideal. From the exact sequence $0\to\mathfrak{a}\to B\to B/\mathfrak{a}\to0$ and
flatness we get $0\to\mathfrak{a}A\to A\to (B/\mathfrak{a})\otimes_B A\to0$, with
$\mathfrak{a}A$ identified with its image $A\mathfrak{a}$ and
$(B/\mathfrak{a})\otimes_BA\cong A/A\mathfrak{a}$. By faithful flatness the natural
map $B/\mathfrak{a}\to (B/\mathfrak{a})\otimes_B A$ is injective: if $0\neq\bar b$ had
zero image, the cyclic module $B/(\mathfrak{a}+Bb)$ would satisfy
$(B/(\mathfrak{a}+Bb))\otimes_B A=0$ and hence would be zero by faithful flatness,
contrary to $\bar b\neq0$. Chasing $\bar b$ through
$B/\mathfrak{a}\to A/A\mathfrak{a}$ gives $B\mathfrak{a}\cap B\subseteq\mathfrak{a}$,
and the other inclusion is trivial.

\emph{(3) Descent.} Let $\mathfrak{a}_1\subseteq\mathfrak{a}_2\subseteq\cdots$ be an
ascending chain of two-sided ideals of $B$. Then $A\mathfrak{a}_1\subseteq
A\mathfrak{a}_2\subseteq\cdots$ is an ascending chain of two-sided ideals of $A$,
hence constant from some $N$ on. For $n\ge N$, step (2) gives
$\mathfrak{a}_n=A\mathfrak{a}_n\cap B=A\mathfrak{a}_N\cap B=\mathfrak{a}_N$.
\end{proof}

\begin{proposition}[Finite codimension: what transfers]\label{prop:codim}
Let $\hl\subseteq\gl$ be a Lie algebra inclusion over $\kk$ with
$d=\dim_\kk(\gl/\hl)<\infty$. Then:
\begin{enumerate}
\item[(i)] $\UU(\gl)$ is a free left $\UU(\hl)$-module and a free right
$\UU(\hl)$-module; a basis is given by the ordered monomials in a basis of any
$\kk$-complement $V$ of $\hl$ in $\gl$.
\item[(ii)] if $\UU(\gl)$ is noetherian, then $\UU(\hl)$ satisfies the ascending chain
condition on two-sided ideals;
\item[(iii)] if in addition $\UU(\gl)$ is finitely generated as a left
$\UU(\hl)$-module, then $\UU(\hl)$ left noetherian implies $\UU(\gl)$ left
noetherian, so the two noetherian properties are equivalent;
\item[(iv)] the extra hypothesis in (iii) is not automatic: see
Example~\ref{ex:nomodule}, which is one-codimensional.
\end{enumerate}
\end{proposition}

\begin{proof}
\emph{(i)} Choose a totally ordered $\kk$-basis of $\gl$ adapted to
$\gl=\hl\oplus V$, all elements of $\hl$ preceding all elements of $V$; this is
possible because $d<\infty$. By PBW the ordered monomials in this basis form a basis
of $\UU(\gl)$, and those monomials whose $\hl$-exponent is empty --- i.e.\ the ordered
monomials in the basis of $V$ --- form a basis of $\UU(\gl)$ as a left
$\UU(\hl)$-module: every ordered monomial in the full basis is a product
$h\,m$ with $h$ an ordered monomial in the $\hl$-part and $m$ an ordered monomial in
the $V$-part, and distinct $m$ give distinct basis elements. The right module
statement is identical, ordering the basis with $V$ first.

\emph{(ii)} By (i) and Theorem~\ref{thm:descent} applied to
$B=\UU(\hl)\subseteq A=\UU(\gl)$.

\emph{(iii)} By (i) and Lemma~\ref{lem:module}.

\emph{(iv)} See Example~\ref{ex:nomodule}, where $d=1$.
\end{proof}

\begin{example}[Finite codimension does not give module finiteness]\label{ex:nomodule}
Let $\gl=\kk x\oplus\kk y\oplus\kk z$ with brackets $[x,z]=z$ and $[x,y]=[y,z]=0$,
and $\hl=\kk x\oplus\kk y$, so that $\dim_\kk\gl=3$, $\dim_\kk\hl=2$ and $d=1$. Then
$\UU(\gl)$ is not finitely generated as a left $\UU(\hl)$-module.
\end{example}

\begin{proof}
\emph{(1) A normal form.} In $\UU(\gl)$ the relation $[x,z]=z$ reads
$zx=xz-z=(x-1)z$; together with $xy=yx$, $zy=yz$ this gives, for every
$f\in\kk[x,y]$, the identity $zf(x,y)=f(x-1,y)z$. Let $\sigma$ be the automorphism
$f\mapsto f(x-1,y)$ of $\kk[x,y]$. Then $zf=\sigma(f)z$, and by induction
$z^nf=\sigma^n(f)z^n$.

\emph{(2) A skew polynomial description.} It follows that every element of
$\UU(\gl)$ has a unique expression $\sum_{n\ge0}f_nz^n$ with $f_n\in\kk[x,y]$ (compare
the PBW basis $x^ay^bz^n$), and that
$(fz^m)(gz^n)=f\sigma^m(g)z^{m+n}$. In particular $\UU(\gl)$ is the skew polynomial
algebra $\kk[x,y][z;\sigma]$.

\emph{(3) Infinite rank.} Left multiplication by $\UU(\hl)=\kk[x,y]$ preserves the
$z$-degree: $f\cdot(gz^n)=(fg)z^n$. Hence, as a left $\kk[x,y]$-module,
$\UU(\gl)=\bigoplus_{n\ge0}\kk[x,y]z^n$, a direct sum of countably infinitely many
nonzero cyclic summands. A finitely generated module is a quotient of a finite direct
sum of copies of $\kk[x,y]$ and therefore has the property that the ascending chain of
submodules $M_0\subseteq M_1\subseteq\cdots$ with $M_j$ the span of the first $j$
summands of an infinite free direct sum cannot be finitely generated (a finite
generating set involves only finitely many $z$-degrees). Concretely, the left
$\kk[x,y]$-submodule $\bigoplus_{n\ge0}\kk[x,y]z^n$ needs the infinitely many
elements $z^n$: no finite set generates it, because the $z$-degree of a combination
$\sum_i f_i m_i$ is at most the maximum $z$-degree of the $m_i$. Hence $\UU(\gl)$ is
not finitely generated over $\UU(\hl)$.
\end{proof}

\begin{remark}\label{rem:codimwhy}
Example~\ref{ex:nomodule} shows that the implication ``finite codimension
$\Rightarrow$ module finiteness'', which would make Proposition~\ref{prop:codim}(iii)
applicable automatically, is false. It fails for a structural reason: commuting an
element of $V$ past an element of $\hl$ produces an element of $\gl$, hence possibly
another $V$-factor, and the $V$-degree is not bounded. Therefore a transfer of
noetherianity from $\UU(\hl)$ to $\UU(\gl)$ needs an additional argument; in the
literature this is supplied case by case, and that argument is not reconstructed in
this note.
\end{remark}

\begin{corollary}[The Witt pair]\label{cor:witt}
Let $W=\Der\kk[t,t^{-1}]$ and $W_{\ge-1}=\Der\kk[t]$, so that $W_{\ge-1}\subset W$ has
codimension $1$. Then:
\begin{enumerate}
\item[(i)] $\UU(W)$ is a free left and right $\UU(W_{\ge-1})$-module;
\item[(ii)] every strictly ascending chain of two-sided ideals of $\UU(W)$ gives one
of $\UU(W_{\ge-1})$, so if $\UU(W)$ is noetherian then $\UU(W_{\ge-1})$ has the
ascending chain condition on two-sided ideals;
\item[(iii)] if the finite-codimension transfer principle holds for this pair (that
is, if left noetherianity transfers in either direction between $\UU(W_{\ge-1})$ and
$\UU(W)$), then the two enveloping algebras are simultaneously noetherian or
simultaneously not; in particular, since $\UU(W)$ is not noetherian
(Corollary~\ref{cor:Wgraded}), $\UU(W_{\ge-1})$ would then not be noetherian either.
\end{enumerate}
\end{corollary}

\begin{proof}
(i) is Proposition~\ref{prop:codim}(i) with $\hl=W_{\ge-1}$, $\gl=W$ and $d=1$:
$W=\bigoplus_{m\ge-1}\kk d_m$ and $W_{\ge-1}=\bigoplus_{m\ge-1}\kk d_m$ inside
$W=\bigoplus_{m\in\mathbb{Z}}\kk d_m$ differ by the one-dimensional space $\kk d_{-1}$.
(ii) is Proposition~\ref{prop:codim}(ii). (iii) is the stated extra hypothesis
together with (ii); the parenthetical non-noetherianity of $\UU(W)$ is
Corollary~\ref{cor:Wgraded} below, which is proved from
Conjecture~\ref{thm:gradedsimple}.
\end{proof}

\subsection{Graded obstructions (reported from the literature)}

\begin{conjecture}[Graded simplicity obstructs noetherianity]\label{thm:gradedsimple}
Let $\kk$ be a field and let $\LL=\bigoplus_{m\in\mathbb{Z}^n}\LL_m$ be a
$\mathbb{Z}^n$-graded Lie algebra over $\kk$ with $[\LL_m,\LL_{m'}]\subseteq
\LL_{m+m'}$. Assume:
\begin{enumerate}
\item[(a)] $\dim_\kk\LL_m<\infty$ for every $m$;
\item[(b)] $\LL$ is graded simple, that is, $\LL$ has no nonzero proper graded ideal;
\item[(c)] $\dim_\kk\LL=\infty$.
\end{enumerate}
Then $\UU(\LL)$ is not noetherian.\footnote{This statement is reported from the
published literature rather than proved in this note: the hypotheses (a)--(c) are
exactly the standing hypotheses on $\mathbb{Z}^n$-graded simple Lie algebras in that
work, and the conclusion is the direction established there. The detailed argument ---
in which the strictness of the chain of truncated ideals is verified inside each
graded-simple family --- is not reconstructed line by line in this manuscript, so the
statement must be read as a report of a known result under its own hypotheses, not as
a self-contained proof. The sketch in the text indicates the shape of the argument and
names the step that is not reproduced here.}
\end{conjecture}

\begin{remark}[Why this is a conjecture and not a theorem here; the missing step]\label{rem:gradedsimple-missing}
\emph{(1) Truncation.} For $r\ge1$ let $J_r\subseteq\UU(\LL)$ be the two-sided ideal
generated by the image of $\bigoplus_{|m|>r}\LL_m$; these form an ascending chain
$J_1\subseteq J_2\subseteq\cdots$ of two-sided ideals.

\emph{(2) Reduction of strictness to the grading.} Because each $\LL_m$ is
finite dimensional and the ideals are generated by homogeneous elements, membership in
$J_r$ can be tested degree by degree; the quotient $\UU(\LL)/J_r$ is a filtered
algebra whose degree-$0$ part ... \emph{(this is the step that is not carried out in
this note)} one must exhibit, for each $r$, a homogeneous element $y_r\in\LL$ of
degree $>r$ which is not in $\LL\cap J_r$, i.e.\ which cannot be rewritten, using the
relations of $\UU(\LL)$ and the generators of $J_r$, into lower degree. In the
literature this is done inside each graded-simple family, using the non-degeneracy of
the brackets among the extreme graded components; graded simplicity (b) is used
precisely to rule out the possibility that the higher graded components generate only
a proper graded ideal.

\emph{(3) Conclusion, granted (2).} If $J_r\subsetneq J_{r+1}$ for all $r$, the chain
in (1) is strictly ascending and $\UU(\LL)$ is not noetherian (an ascending chain
condition on two-sided ideals fails, hence left noetherianity fails).
\end{remark}

\begin{corollary}[The Witt algebra]\label{cor:Wgraded}
$W=\Der\kk[t,t^{-1}]$ is $\mathbb{Z}$-graded with $\dim_\kk W_m=1$ for all
$m\in\mathbb{Z}$ and is graded simple; hence $\UU(W)$ is not noetherian.
\end{corollary}

\begin{proof}
The grading is $W=\bigoplus_{m\in\mathbb{Z}}\kk d_m$ with $d_m=t^{m+1}\partial_t$ and
$[d_m,d_n]=(n-m)d_{m+n}$, so $\dim_\kk W_m=1<\infty$ and
$[W_m,W_{m'}]\subseteq W_{m+m'}$. Graded simplicity is the standard fact that a
nonzero graded ideal contains some $d_m$ and then, because
$[d_{m+1},d_{-1}]=(m+2)d_m$ and $(n-m)\neq0$ in characteristic $0$ for $n\neq m$, all
brackets generated by translates of $d_m$ span $W$; we record it here without proof as
part of the standard structure theory of $W$. Applying
Conjecture~\ref{thm:gradedsimple} with $n=1$, (a) and (b) are satisfied, and (c) holds
because $\dim_\kk W=\infty$. Hence $\UU(W)$ is not noetherian.
\end{proof}

\begin{conjecture}[Krichever--Novikov algebras]\label{thm:kn}
Let $X$ be a smooth projective curve of genus $g$ over $\kk$ with two marked points
$P_+,P_-$, and let $L(X,P_+,P_-)$ be the associated Krichever--Novikov Lie algebra, the
Lie algebra of meromorphic vector fields on $X$ holomorphic away from $P_\pm$ and with
poles of bounded order at the marked points. Under the standard hypotheses
(fixed genus and fixed orders of poles at the two marked points, so that the algebra is
finitely generated over the ring of functions on $X$) one has: $L(X,P_+,P_-)$ is
$\mathbb{Z}$-graded with finite-dimensional homogeneous components, it is simple for
the standard choices of the generating fields, and it is infinite dimensional.
Under these hypotheses $\UU(L(X,P_+,P_-))$ is not noetherian, in the direction of the
known results; the case $g=0$ with the standard one-point presentation is the Witt
algebra of Corollary~\ref{cor:Wgraded}, while $g\ge1$ is the genuinely new case.\footnote{As with
Conjecture~\ref{thm:gradedsimple}: the hypotheses listed are the ones under which the
result is published, and the conclusion is reported in the direction established
there. The proof is not reconstructed here; in particular we do not verify the
graded-simple hypothesis for $g\ge1$ line by line, which is why it is phrased as an
assumption in the statement. The statement is asserted only under these hypotheses and
says nothing about Krichever--Novikov algebras for other choices of the point data.}
\end{conjecture}

\begin{remark}[Why this is a conjecture here; the missing step: strictness of the truncated chain inside each graded-simple family]\label{rem:kn-missing}
\emph{(1) Grading and finiteness.} By the definition of
$L(X,P_+,P_-)$ as the Lie algebra of meromorphic vector fields on $X$ that are
holomorphic away from $P_\pm$ and have poles of bounded order there, the order of the
pole at $P_+$ (equivalently: minus the order at $P_-$, the two being linked by the
residue) defines a $\mathbb{Z}$-grading $L(X,P_+,P_-)=\bigoplus_{m\in\mathbb{Z}}L_m$,
and $[L_m,L_{m'}]\subseteq L_{m+m'}$; with the genus and the pole orders fixed the space
of fields with a given pole bound at both points is finite dimensional, so
$\dim_\kk L_m<\infty$ for every $m$. The algebra is infinite dimensional because the
pole order is unbounded.

\emph{(2) Graded simplicity, as an assumption of the statement.} For the standard
choices of the generating fields the algebra is graded simple; this is one of the
hypotheses listed in the statement and is not verified here line by line. Graded
simplicity is used exactly as hypothesis (b) of the reported graded obstruction.

\emph{(3) Conclusion.} Under (1) and (2) hypotheses (a), (b), (c) of the graded
obstruction are satisfied with $n=1$, and the reported obstruction gives that
$\UU(L(X,P_+,P_-))$ is not noetherian. The case $g=0$ is the Witt algebra
(Corollary~\ref{cor:Wgraded}); the case $g\ge1$ is the genuinely new case.
\end{remark}

\begin{conjecture}[Graded algebras of polynomial growth]\label{thm:poly}
Let $\LL=\bigoplus_{m\in\mathbb{Z}}\LL_m$ be a $\mathbb{Z}$-graded Lie algebra over
$\kk$ with $\dim_\kk\LL_m<\infty$ for all $m$ and $\dim_\kk\LL=\infty$. Assume that the
growth of the homogeneous components is polynomial, i.e.\ there are $C>0$ and $q\ge0$
with $\sum_{|m|\le r}\dim_\kk\LL_m\le Cr^{q}$ for all $r\ge1$, and that $\LL$ is
graded simple. Then $\UU(\LL)$ is not noetherian. More generally, the conclusion is
claimed in the literature for graded algebras in that class, and only there.\footnote{Reported, not reconstructed: the polynomial-growth hypothesis replaces the
uniform boundedness of the homogeneous components used in
Conjecture~\ref{thm:gradedsimple}, and the argument for the strictness of the truncated
chain must then control the growth degree by degree. This note does not reproduce that
control, so the statement is asserted under the hypotheses as stated and no others.}
\end{conjecture}

\begin{remark}[Why this is a conjecture here; the missing step: the degree-by-degree growth control for the truncated chain]\label{rem:poly-missing}
\emph{(1) Grading and finiteness.} By hypothesis $\LL=\bigoplus_{m\in\mathbb{Z}}\LL_m$
with $[\LL_m,\LL_{m'}]\subseteq\LL_{m+m'}$ and $\dim_\kk\LL_m<\infty$ for every $m$.
Every element of $\LL$ therefore lies in a finite sum of finitely many homogeneous
components, each of which is a finite-dimensional $\kk$-vector space; this is a
restatement of the hypotheses, not an independent result of this note.

\emph{(2) The growth control (the step not reconstructed here).} The polynomial-growth
hypothesis $\sum_{|m|\le r}\dim_\kk\LL_m\le Cr^{q}$ is what replaces the uniform
boundedness of the homogeneous components used in
Conjecture~\ref{thm:gradedsimple}: for the truncated chain to be shown strict one must
compare the graded components degree by degree against the growth function. That
comparison is not carried out in this note.

\emph{(3) Conclusion, conditionally on the imported direction.} Granting the published
argument for graded algebras of polynomial growth, (1) and (2) put $\LL$ inside its
hypotheses and $\UU(\LL)$ is not noetherian. The step is an import from the literature,
not a proof produced here, and the statement holds only under the hypotheses written
above and under no others.
\end{remark}

\begin{conjecture}[Finitely generated infinite-step solvable algebras]\label{thm:solvable}
Let $N\ge1$, let $S_N$ denote the class of finitely generated infinite-dimensional
Lie algebras $\gl$ over $\kk$ that are solvable of derived length $N$, i.e.\
$D_N\gl=0\neq D_{N-1}\gl$ for the derived series $D_0\gl=\gl$,
$D_{k+1}\gl=[D_k\gl,D_k\gl]$. Then, in the direction of the known results, for every
$\gl\in S_N$ with $N\ge1$ the enveloping algebra $\UU(\gl)$ is not noetherian, with one
exception: in the class $S_2$ there are algebras whose enveloping algebra fails to be
noetherian while being a Krull domain, so that non-noetherianity there cannot be
detected by the failure of being a domain or of being a Krull domain.\footnote{Reported, not reconstructed. The statement is asserted only for the class $S_N$
as defined (finitely generated, infinite dimensional, solvable of derived length $N$)
and says nothing about solvable algebras that are not finitely generated, nor about the
converse direction. In particular the ``$S_2$ but Krull'' clause is a report from the
literature and is not verified here; it is included because it delimits which
ring-theoretic properties can be used as a test for noetherianity in the solvable
class.}
\end{conjecture}

\begin{remark}[Why this is a conjecture here; the missing step: an explicit strictly ascending chain of two-sided ideals, with its strictness certified inside each derived layer]\label{rem:solvable-missing}
\emph{(1) Reading the class.} By definition $D_0\gl=\gl$ and
$D_{k+1}\gl=[D_k\gl,D_k\gl]$, so $\gl\in S_N$ means: $\gl$ is finitely generated,
$\dim_\kk\gl=\infty$, and $D_N\gl=0\neq D_{N-1}\gl$. Thus $\gl$ has a finite
generating set $x_1,\dots,x_r$ and a finite derived series of length $N$; both facts
are hypotheses here, not consequences derived in this note.

\emph{(2) The chain step (imported, not reconstructed).} What has to be shown for a
member of $S_N$ is the existence of a strictly ascending chain of two-sided ideals of
$\UU(\gl)$. Since $\gl$ is infinite dimensional while $D_N\gl=0$, one filters $\gl$ by
the ideals of the derived series and passes to the induced filtration of $\UU(\gl)$;
the strictness of the resulting chain must be verified inside each derived layer using
finite generation of $\gl$ and the non-vanishing of the last nonzero derived ideal
$D_{N-1}\gl$. That verification is the step of the published argument, and it is
\emph{not} carried out here: no chain is written down explicitly in this note, and no
rank computation certifying its strictness is performed. Consequently the conclusion
below is a conditional one, valid only in the direction established in the literature
under the hypotheses of the statement.

\emph{(3) Conclusion, granted (2).} If such a strictly ascending chain of two-sided
ideals exists, then $\UU(\gl)$ fails the ascending chain condition on two-sided
ideals, hence fails to be left noetherian (a left noetherian ring satisfies the
ascending chain condition on two-sided ideals as well). The parenthetical clause about
$S_2$ and Krull domains is likewise a report from the literature and is not checked
here; it is recorded only to delimit which ring-theoretic properties can serve as a
test for noetherianity in this class.
\end{remark}

\section{Discussion and open problems}\label{sec:open}

The positive direction is settled (Theorem~\ref{thm:findim}), the abelian case is
settled (Theorem~\ref{thm:abelian}), and the graded obstruction covers the graded
simple families with finite-dimensional homogeneous components
(Conjecture~\ref{thm:gradedsimple} and its instances). What remains is exactly the middle
ground.

\begin{conjecture}[Sierra--Walton, as formulated in \cite{SW14}]\label{conj:SW}
Let $\kk$ be a field of characteristic $0$ and let $\gl$ be a Lie algebra over $\kk$.
If $\dim_\kk\gl=\infty$, then $\UU(\gl)$ is not noetherian. Equivalently: $\UU(\gl)$
noetherian implies $\dim_\kk\gl<\infty$.
\end{conjecture}

\begin{openproblem}[What is missing]\label{op:missing}
Conjecture~\ref{conj:SW} is open. Concretely, the following are missing, and each is a
separate task:
\begin{enumerate}
\item[(1)] \emph{A construction.} No example is known of an infinite-dimensional
non-abelian $\gl$ with $\UU(\gl)$ noetherian, and no example is known of an
infinite-dimensional $\gl$ with $\UU(\gl)$ noetherian other than in the abelian
(hence non-noetherian) direction. What is missing is either such a construction, or a
theorem excluding it under stated extra hypotheses.
\item[(2)] \emph{A converse under additional hypotheses.} There is no theorem of the
form ``if $\gl$ is infinite dimensional and nilpotent (resp.\ solvable, resp.\ of
bounded growth, resp.\ graded in some class other than the graded-simple ones above),
then $\UU(\gl)$ is not noetherian'', except in the classes listed in
Conjectures~\ref{thm:kn}, \ref{thm:poly} and \ref{thm:solvable}, each of which holds only
under its own hypotheses.
\item[(3)] \emph{A transfer principle for finite-codimension subalgebras.} The step
``$\UU(\hl)$ noetherian $\Leftrightarrow$ $\UU(\gl)$ noetherian for
$\dim_\kk(\gl/\hl)<\infty$'' is used for individual pairs in the literature, but
Proposition~\ref{prop:codim} shows that the naive route through module finiteness is
not available in general (Example~\ref{ex:nomodule}); a proof of the transfer
principle, or a counterexample to it, is not in this note.
\item[(4)] \emph{Independence of the classes.} The graded obstruction is stated per
graded-simple family; it is not known here how far the several families overlap, and
no claim is made that the list in Section~3 is exhaustive.
\end{enumerate}
\end{openproblem}

\begin{remark}[On the status of the statements]
Items (1)--(4) are the honest gaps. The results of Section~3 that are proved here
(Theorem~\ref{thm:findim}, Theorem~\ref{thm:abelian}, Lemma~\ref{lem:pbw},
Lemma~\ref{lem:hilbert}, Lemma~\ref{lem:lift}, Lemma~\ref{lem:module},
Theorem~\ref{thm:descent}, Proposition~\ref{prop:codim},
Example~\ref{ex:nomodule}) are proved under exactly the hypotheses stated, in
characteristic $0$ where the anti-automorphism argument is used. The results reported
rather than proved (Conjecture~\ref{thm:gradedsimple} and its instances
\ref{cor:Wgraded}, \ref{thm:kn}, \ref{thm:poly}, \ref{thm:solvable}) are asserted only
under their published hypotheses and are marked in place; the Witt case
(Corollary~\ref{cor:Wgraded}) rests on the standard structure theory of $W$, which is
also not reconstructed here.
\end{remark}

\begin{thebibliography}{9}
\bibitem[KN87]{KN87} I.~M. Krichever and S.~P. Novikov,
\emph{Algebras of Virasoro type, Riemann surfaces and structures of the theory of
solitons}, Funktsional. Anal. i Prilozhen. \textbf{21} (1987), no.~2, 46--63.
\bibitem[MR01]{MR01} J.~C. McConnell and J.~C. Robson,
\emph{Noncommutative Noetherian Rings}, Graduate Studies in Mathematics \textbf{30},
American Mathematical Society, Providence, RI, 2001.
\bibitem[SW14]{SW14} S.~Sierra and C.~Walton,
\emph{The universal enveloping algebra of the Witt algebra is not noetherian},
Adv. Math. \textbf{262} (2014), 239--260.
\end{thebibliography}

\appendix
\section{Computational evidence}\label{app:comp}

This appendix collects all computational evidence referred to anywhere in this paper:
the two finite computations (A) and (B), the algorithm used for each, its parameter
domain, the resulting quantities, and the internal record identifiers together with how
to rerun them. Nothing here enters the proofs of Section~3; every conclusion that
rests on these computations depends on that evidence and on the stated bounds, and no
computation below is a proof of anything.

\paragraph{A.1 Computation (A): filtered dimensions in dimension three.}
\emph{Quantity computed.} The dimension of the $n$-th filtered piece $\UU_n(\gl)$ for
$\dim_\kk\gl=3$ and $n=0,1,\dots,8$, and the third finite difference of that sequence.
\emph{Algorithm.} Ordered monomials in a basis of $\gl$ are multiplied modulo the span
of the monomials of lower total degree, and the dimension is read off as a rank over
$\kk$; this is a direct comparison with the PBW basis of Lemma~\ref{lem:pbw}.
\emph{Parameter domain.} $d=3$, with the abelian, Heisenberg and $\mathfrak{sl}_2$
brackets, and $0\le n\le8$; no other dimension and no larger $n$ was run.

\begin{table}[h]
\centering
\begin{tabular}{rrr}
\toprule
$n$ & $\dim\UU_n$ (all three algebras) & third finite difference \\
\midrule
0 & 1 & -- \\
1 & 4 & -- \\
2 & 10 & -- \\
3 & 20 & 1 \\
4 & 35 & 1 \\
5 & 56 & 1 \\
6 & 84 & 1 \\
7 & 120 & 1 \\
8 & 165 & 1 \\
\bottomrule
\end{tabular}
\caption{Computation (A): the computed dimension agrees with $\binom{n+3}{3}$ for
$n\le8$, and the third finite difference is $1$ on the computed range. Bound of the
computation: $\dim\gl=3$, $n\le8$.}
\end{table}

\paragraph{A.2 Computation (B): strictness of a truncated family of ideals.}
\emph{Quantity computed.} For $n=1,\dots,12$ the inclusion $I_n\subsetneq I_{n+1}$ for
$I_n=(x_1,\dots,x_n)$, inside a polynomial ring in countably many indeterminates
truncated at degree $12$.
\emph{Algorithm.} For each $n$ a separating functional on the degree-$(n{+}1)$ part is
exhibited exactly: the algebra map killing $x_1,\dots,x_n$ and fixing $x_{n+1}$
vanishes on $I_n$ and is nonzero on $x_{n+1}$. The degree truncation enters only the
bookkeeping of which monomials were enumerated, not the certificate.
\emph{Parameter domain.} $1\le n\le12$, degree bound $12$; no larger $n$ was run.

\begin{table}[h]
\centering
\begin{tabular}{rrrl}
\toprule
family index $n$ & degree bound & enumerated monomials (order) & inclusion strict? \\
\midrule
1--12 & 12 & $10^{2}$--$10^{3}$ & yes, certified exactly \\
\bottomrule
\end{tabular}
\caption{Computation (B): for each $n\le12$ the inclusion $I_n\subsetneq I_{n+1}$ is
certified by a separating functional on the degree-$(n{+}1)$ part. Bound: $n\le12$.}
\end{table}

\paragraph{Computations.}
\begin{itemize}
\item \texttt{exp:000011} --- filtered dimension count, $\dim_\kk\gl=3$ (abelian,
Heisenberg, $\mathfrak{sl}_2$), $n=0,\dots,8$; parameter set
$\{d=3,\ n_{\max}=8\}$; script digest recorded with the experiment
(\texttt{e435e239ba21f72f}); conclusion recorded as \texttt{supports} with the explicit
scope note ``$d=3$ only; does not address noetherianity''. Rerun: same script, same
parameters.
\item \texttt{exp:000012} --- truncated family of ideals $I_n$, $n=1,\dots,12$,
members $I_1,\dots,I_{12}$, exact separation test on the degree-$(n{+}1)$ part;
parameter set $\{n_{\max}=12\}$; script digest recorded with the experiment
(\texttt{3473d47b6f460c3a}); conclusion recorded as \texttt{supports}.
\item Both records have \texttt{method=numeric} and are the source of the two tables in
Appendix~\ref{app:comp}. Rerunning either with a different $n_{\max}$ or a different
$d$ requires a new record; the present tables must not be extrapolated.
\end{itemize}

\paragraph{Statements used.}
\begin{itemize}
\item \texttt{stmt:000007} --- finite dimension implies noetherian (statement level,
no proof text locally).
\item \texttt{stmt:000004} --- infinite-dimensional abelian implies not noetherian.
\item \texttt{stmt:000005} --- finite-codimension transfer, in the form discussed in
Proposition~\ref{prop:codim}.
\item \texttt{stmt:000010} --- graded-simple infinite-dimensional implies not
noetherian (literature direction), the source of Conjecture~\ref{thm:gradedsimple}.
\item \texttt{stmt:000018}, \texttt{stmt:000019} --- derived series and polynomial
growth definitions used in Conjectures~\ref{thm:solvable} and \ref{thm:poly}.
\item \texttt{stmt:000001}, \texttt{stmt:000002} --- Krichever--Novikov algebras and
the two conventions for the Witt algebra (used in Definition~\ref{def:witt}).
\item \texttt{stmt:000003} --- the graded-simple definition and the
finite-dimensional-homogeneous-components convention.
\end{itemize}

\paragraph{Literature items.}
\begin{itemize}
\item \texttt{lit:000001} --- source for the Krichever--Novikov and Witt definitions
(quotation verified verbatim at the pages recorded there).
\item \texttt{lit:000002} --- source for the graded-simple setting and the two Witt
conventions (quotation verified verbatim).
\item \texttt{lit:000004} --- McConnell--Robson, \emph{Noncommutative Noetherian
Rings}: bibliographic record only, no local full text, so nothing in this note is
quoted from it; the reference to Corollary~1.7.4 for the positive direction is
recorded at \texttt{stmt:000012}.
\item \texttt{lit:000024} --- source for the derived-series and polynomial-growth
definitions (quotation verified verbatim).
\item \texttt{lit:000025} --- bibliographic record only (no local full text); it is
\emph{not} a source of any verbatim quotation in this note.
\end{itemize}

\paragraph{How to reproduce the two computations.} Both are deterministic scripts run
in the workspace; each rerun must pass the parameters named above and must be compared
against the digests recorded with \texttt{exp:000011} and \texttt{exp:000012}. A rerun
with a different parameter set produces a different record and does not extend the
bounds stated in Appendix~\ref{app:comp}.

\end{document}
